\section{Permisos y SFTP}

La raíz de WordPress pertenece a \texttt{wpadmin:www-data}. Los directorios usan modo 750 y los archivos modo 640. El propietario puede modificar contenido; Apache, miembro del grupo \texttt{www-data}, puede leer y atravesar directorios; los demás usuarios no reciben acceso.

\subsection{Aplicación reproducible de \texttt{chown} y \texttt{chmod}}

Los permisos se aplicaron con el script revisado \texttt{scripts/aplicar-permisos-wordpress.sh}. Primero se ejecutó sin \texttt{--apply} para validar la ruta y los usuarios; después se aplicaron las operaciones equivalentes del Listado~\ref{lst:permisos-wordpress}. Se evitó \texttt{chmod 777} porque otorgaría privilegios innecesarios a terceros.

\begin{lstlisting}[caption={Comandos aplicados a la carpeta de WordPress.},label={lst:permisos-wordpress}]
sudo chown -R wpadmin:www-data /var/www/wordpress
sudo find /var/www/wordpress -type d -exec chmod 750 {} \;
sudo find /var/www/wordpress -type f -exec chmod 640 {} \;
\end{lstlisting}

\texttt{chown} establece a \texttt{wpadmin} como propietario y a \texttt{www-data} como grupo. En directorios, 750 concede lectura, escritura y recorrido al propietario; lectura y recorrido al grupo; ningún permiso a otros. En archivos, 640 concede lectura y escritura al propietario, lectura al grupo y ningún permiso a otros. El resultado esperado era que Apache pudiera servir el sitio como miembro de \texttt{www-data}, que \texttt{wpadmin} pudiera modificarlo por SFTP y que un usuario ajeno no pudiera escribir.

La comprobación posterior utilizó \texttt{stat -c '\%U \%G \%a \%n'} sobre la raíz, un archivo y un directorio. La prueba negativa intentó crear un archivo como \texttt{wpviewer} y obtuvo \textit{Permission denied}. Así, la gestión de permisos se relaciona directamente con el objetivo de comprender propietario, grupo y otros, y no se limita a describir valores sin probar su efecto.

\begin{table}[H]
  \centering
  \begin{tabular}{llll}
    \toprule
    \textbf{Objeto} & \textbf{Propietario} & \textbf{Grupo} & \textbf{Modo} \\
    \midrule
    \texttt{/var/www/wordpress} & \texttt{wpadmin} & \texttt{www-data} & 750 \\
    Directorios descendientes & \texttt{wpadmin} & \texttt{www-data} & 750 \\
    Archivos WordPress & \texttt{wpadmin} & \texttt{www-data} & 640 \\
    \bottomrule
  \end{tabular}
  \caption{Modelo de propiedad y permisos aplicado.}
\end{table}

FileZilla se configuró con protocolo SFTP, host \texttt{10.50.5.20}, puerto 22222, usuario \texttt{wpadmin} y llave dedicada \texttt{lab5\_wpadmin}. No se utilizaron FTP ni credenciales de contraseña. La prueba positiva accedió al árbol de WordPress; la negativa utilizó un usuario sin autorización de escritura.

\evidencia{FileZilla conectado mediante SFTP}{05-seguridad/E12-filezilla-sftp.png}{La interfaz muestra la sesión SFTP del usuario \texttt{wpadmin} y el contenido remoto de WordPress. Esto demuestra que la administración de archivos parte del cliente gráfico.}

\evidencia{Escritura denegada al usuario no autorizado}{05-seguridad/E13-permiso-denegado.png}{El intento del usuario de prueba falla por permisos. La prueba negativa verifica que el acceso general al servidor no equivale a permiso de modificación sobre WordPress.}
